\documentclass[a4paper,12pt]{extarticle}
\usepackage{geometry}
\usepackage[T2A,T1]{fontenc}% T2A до babel — кириллица; T1 — латиница (см. предупреждение babel)
\usepackage[utf8]{inputenc}
\usepackage[english,russian]{babel}
\usepackage{amsmath}
\usepackage{amsthm}
\usepackage{amssymb}
\usepackage{fancyhdr}
\usepackage{setspace}
\usepackage{graphicx}
\usepackage{colortbl}
\usepackage{tikz}
\usepackage{pgf}
\usepackage{subcaption}
\usepackage{listings}
\lstset{breaklines=true,columns=fullflexible}% длинные строки в lstlisting не уезжают вправо
\usepackage{indentfirst}
\usepackage[
backend=biber,
style=numeric,
maxbibnames=99
]{biblatex}
\addbibresource{refs.bib}
\usepackage[colorlinks,citecolor=blue,linkcolor=blue,bookmarks=false,hypertexnames=true, urlcolor=blue]{hyperref}
\usepackage{xurl}% перенос длинных путей и имён (в отличие от \texttt)
\makeatletter
\g@addto@macro\UrlBreaks{\do\/\do\_\do\.\do\-\do\:\do\=\do\?\do\&\do\#\do\<\do\>}
\makeatother
\urlstyle{tt}
\newcommand{\ppath}[1]{{\hypersetup{urlcolor=black}\path{#1}}}% путь/идентификатор, перенос по / и _
\newcommand{\httpapi}[1]{{\hypersetup{urlcolor=black}\path|#1|}}% HTTP-метод и путь (есть пробел)
\setlength{\emergencystretch}{4em}% последний рубеж против переполнения абзаца
\usepackage{mathtools}
\usepackage{booktabs}
\usepackage[flushleft]{threeparttable}
\usepackage{tablefootnote}

\usepackage{chngcntr} % нумерация графиков и таблиц по секциям
\counterwithin{table}{section}
\counterwithin{figure}{section}

\graphicspath{{graphics/}}%путь к рисункам

\makeatletter
% \renewcommand{\@biblabel}[1]{#1.} % Заменяем библиографию с квадратных скобок на точку:
\makeatother

\geometry{left=2.5cm}% левое поле
\geometry{right=1.0cm}% правое поле
\geometry{top=2.0cm}% верхнее поле
\geometry{bottom=2.0cm}% нижнее поле
\setlength{\parindent}{1.25cm}
\renewcommand{\baselinestretch}{1.5} % междустрочный интервал


\newcommand{\bibref}[3]{\hyperlink{#1}{#2 (#3)}} % biblabel, authors, year
\addto\captionsrussian{\def\refname{Список литературы (или источников)}}

\renewcommand{\theenumi}{\arabic{enumi}}% Меняем везде перечисления на цифра.цифра
\renewcommand{\labelenumi}{\arabic{enumi}}% Меняем везде перечисления на цифра.цифра
\renewcommand{\theenumii}{.\arabic{enumii}}% Меняем везде перечисления на цифра.цифра
\renewcommand{\labelenumii}{\arabic{enumi}.\arabic{enumii}.}% Меняем везде перечисления на цифра.цифра
\renewcommand{\theenumiii}{.\arabic{enumiii}}% Меняем везде перечисления на цифра.цифра
\renewcommand{\labelenumiii}{\arabic{enumi}.\arabic{enumii}.\arabic{enumiii}.}% Меняем везде перечисления на цифра.цифра

\begin{document}
\input{title_kr}% это титульный лист - выберите подходящий вам из имеющихся в проекте вариантов (kr - курсовая работа у 3 курса, vkr - выпускная квалификационная работа у 4 курса)
\newpage
\setcounter{page}{2}

{
	\hypersetup{linkcolor=black}
	\tableofcontents
}

\newpage
\section*{Аннотация}

Данная работа описывает разработанную систему усиленной неквалифицированной электронной подписи (НЭП) для корпоративной платформы документооборота в соответствии с требованиями Федерального закона № 63-ФЗ «Об электронной подписи»~\cite{fz63} и статьи 22.3 Трудового кодекса РФ~\cite{tk_rf}.
Целью было обеспечить юридически значимый внутренний электронный документооборот, основным преимуществом которого является полная независимость от сторонних удостоверяющих центров и коммерческих средств криптографической защиты информации (СКЗИ), что позволяет значительно снизить эксплуатационные расходы и упрощает интеграцию в существующую корпоративную инфраструктуру.

Реализована микросервисная архитектура. Разработан HTTP API на C++ в фреймворке userver~\cite{userver} с доступом по Bearer-токену, реализованы генерация пары RSA-2048~\cite{rivest1978,rfc3447} средствами OpenSSL~\cite{openssl}, экспорт приватного ключа в PEM с шифрованием AES-256-CBC от пароля подписи и вычисление SHA-256-отпечатка публичного ключа. Криптографические операции выполняются при помощи отдельного микросервиса на Python с использованием библиотек cryptography~\cite{cryptography_python} и pyHanko~\cite{pyhanko}, формирующего отсоединённые подписи PKCS\#7 (CMS)~\cite{pkcs7,rfc5652} для PDF и выполняющего проверку с раздельными признаками \ppath{integrity_ok} и \ppath{signature_ok}. Ключи и метаданные подписей хранятся в PostgreSQL~\cite{postgresql}, документы и подписи — в S3-совместимом объектном хранилище.

Система предоставляет RESTful API~\cite{rest_api}: базовый контур НЭП (генерация ключей, подписание, проверка), а также сценарии с несколькими подписями одного документа, загрузкой КЭП с клиента и выгрузкой ZIP-архива «документ + все подписи». Контракты отражены в спецификации OpenAPI проекта. Корректность подтверждается интеграционными тестами API и юнит-тестами модуля подписи. Для оценки задержек HTTP-ручек предусмотрен сценарий \ppath{scripts/bench_nep_handlers.sh}.

Разработанная система гарантирует целостность и неизменность подписанных документов, идентификацию подписанта и невозможность отказа от авторства подписи.


\addcontentsline{toc}{section}{Аннотация}

\section*{Ключевые слова}
Электронная цифровая подпись, криптография, RSA, SHA-256, PKCS\#7, C++, Python, OpenSSL, корпоративный документооборот, микросервисная архитектура

\pagebreak

\section{Введение}

\subsection{Предметная область}

В современных организациях электронный документооборот становится стандартом деловой практики: договоры, приказы, согласования и кадровые документы всё чаще существуют только в цифровом виде. Чтобы такие документы имели юридическую значимость, недостаточно просто сохранить файл на диске. Нужны средства, которые позволяют убедиться, что текст не изменялся после согласования, идентифицируют авторов подписей и не позволяют им позже отказаться от авторства.

В российской правовой модели эту роль играют электронные подписи. Их виды и условия применения задаёт Федеральный закон № 63-ФЗ «Об электронной подписи»~\cite{fz63}. Для внутреннего корпоративного документооборота наиболее подходящим решением является усиленная неквалифицированная электронная подпись, которая обеспечивает криптографическую защиту без необходимости обращения к аккредитованным удостоверяющим центрам.

Статья 22.3 Трудового кодекса РФ~\cite{tk_rf} легализует использование электронного документооборота в трудовых отношениях при условии соблюдения требований к идентификации сторон и обеспечению целостности документов. Это открывает возможности для цифровизации кадрового делопроизводства и внутренних согласований.


\subsection{Постановка задачи}

Целью данной работы является разработка безопасного и производительного бэкенд-сервиса на языке C++ для создания и проверки электронных цифровых подписей документов в рамках корпоративной платформы. Система должна обеспечивать полный цикл работы с электронной подписью: генерацию ключевых пар для сотрудников, подписание документов и проверку подлинности подписей.

Разрабатываемая система должна соответствовать следующим требованиям:

\begin{enumerate}
    \item \textbf{Криптографическая стойкость.} Использование современных криптографических алгоритмов с достаточным уровнем стойкости: RSA-2048~\cite{rivest1978,rfc3447} для асимметричного шифрования, SHA-256~\cite{sha256} для хеширования.

    \item \textbf{Соответствие стандартам.} Формат подписей должен соответствовать международному стандарту PKCS\#7 (CMS)~\cite{pkcs7,rfc5652}, что обеспечит совместимость с другими системами.

    \item \textbf{Производительность.} Система должна обрабатывать большое количество операций подписания и проверки с минимальными задержками, что критично для корпоративной платформы с тысячами пользователей.

    \item \textbf{Безопасность.} Приватные ключи сотрудников должны надежно храниться, доступ к операциям подписания должен осуществляться только после аутентификации.

    \item \textbf{Аудит.} Все операции с электронными подписями должны логироваться.

    \item \textbf{Интеграция.} Система должна предоставлять удобный API для интеграции с другими компонентами корпоративной платформы.
\end{enumerate}

\subsection{Актуальность и значимость}

Разработка собственной системы электронной подписи для корпоративного документооборота актуальна по следующим причинам:
\begin{enumerate}
    \item Использование коммерческих средств криптографической защиты информации (СКЗИ) и сертифицированных удостоверяющих центров требует значительных финансовых затрат. Для внутреннего документооборота разработка собственного решения позволяет существенно снизить эксплуатационные расходы;
    \item Коммерческие решения часто имеют ограничения в интеграции с существующими корпоративными системами. Собственная разработка позволяет создать API, максимально соответствующий архитектуре платформы;
    \item При использовании внешних сервисов организация зависит от их доступности и политики предоставления услуг. Собственное решение обеспечивает полный контроль над процессами генерации ключей, подписания и проверки документов;
    \item Федеральный закон 63-ФЗ~\cite{fz63} не требует обязательного использования квалифицированной электронной подписи для внутрикорпоративного документооборота. Разработанная система будет создавать усиленные неквалифицированные электронные подписи (НЭП) в виде отсоединенных файлов формата PKCS\#7~\cite{pkcs7}, что обеспечивает гибкость архитектуры и позволяет платформе также работать с квалифицированными электронными подписями (КЭП), получаемыми от клиентских приложений.
\end{enumerate}


\subsection{Достигнутые результаты работы}

Спроектирована микросервисная архитектура с разделением ответственности между компонентами: C++ бэкенд для обработки API-запросов и управления данными, генерации ключевых пар сотрудников, Python микросервис для формирования подписей, проверки их подлинности, PostgreSQL база данных для хранения ключей и метаданных.

Публично доступны три HTTP-эндпоинта НЭП; дополнительно --- загрузка КЭП, цепочка согласования, выгрузка архива подписей (ZIP). Контракт описан в OpenAPI репозитория.

Проведено функциональное тестирование корректности криптографических операций и способности системы обнаруживать модификации подписанных документов. Корректность сценариев подтверждается интеграционными тестами API и юнит-тестами Python-сервиса. Для оценки задержек --- скрипт \ppath{scripts/bench_nep_handlers.sh}.


\subsection{Новизна и значимость результатов}

Основная новизна работы заключается в создании интегрированного решения для электронной подписи, оптимизированного для конкретной корпоративной платформы. В отличие от универсальных коммерческих решений, разрабатываемая система учитывает специфику существующей архитектуры и требования к производительности.

Практическая значимость работы состоит в возможности применения разработанной системы для цифровизации документооборота в организации: подписания внутренних приказов и распоряжений, согласования документов между подразделениями, подтверждения ознакомления сотрудников с документами. Это позволяет сократить время обработки документов, обеспечивает их юридическую значимость и снижает затраты на бумажный документооборот.

\subsection{Структура работы}

Дальнейшая структура работы включает следующие разделы:

\begin{itemize}
    \item \textbf{Раздел 2.} Обзор литературы: криптографические основы, стандарты PKCS\#7/CMS и PKCS\#1, нормативная база, сравнение с коммерческими и облачными решениями, обоснование микросервисного подхода (C++/Python), позиционирование работы.

    \item \textbf{Раздел 3.} Реализация: архитектура НЭП, данные в PostgreSQL и S3, публичный и внутренний HTTP API, множественные подписи и КЭП с клиента, цепочка согласования, ZIP-архив при скачивании, пользовательский сценарий, типовые отказы, объёмы кода, сборка и миграции.

    \item \textbf{Раздел 4.} Криптографические модули на C++ (OpenSSL) и Python (\texttt{NEPSigner}).

    \item \textbf{Раздел 5.} Тестирование (57 автотестов), сценарий замера задержек, сопоставление с аналогами.

    \item \textbf{Раздел 6.} Заключение; список литературы; приложения (сокращения; пример применения миграции БД).
\end{itemize}

\section{Обзор литературы}

\subsection{Теоретические основы криптографии}

Теоретические основы современной криптографии с асимметричным шифрованием были заложены в работе Rivest, Shamir и Adleman~\cite{rivest1978}, которые предложили алгоритм RSA. Этот алгоритм основан на вычислительной сложности факторизации больших целых чисел и до сих пор остается одним из наиболее распространенных методов для создания электронных подписей. В работе использовался RSA с длиной ключа 2048 бит, что в соответствии с современными рекомендациями~\cite{schneier,menezes} обеспечивает достаточный уровень криптографической стойкости.

Для хеширования документов используется алгоритм SHA-256~\cite{sha256}, входящий в семейство Secure Hash Algorithm, стандартизированный NIST. SHA-256 обеспечивает криптографически стойкую однонаправленную функцию хеширования, необходимую для создания цифровой подписи. Важным свойством SHA-256 является отсутствие известных практических атак на поиск коллизий, что критично для обеспечения целостности подписанных документов.

Комплексное изложение теоретических основ криптографии представлено в классических работах Schneier~\cite{schneier} и Menezes et al.~\cite{menezes}, которые подробно описывают математические основы криптографических алгоритмов, их свойства и области применения.

\subsection{Стандарты электронной подписи}

Формат электронной подписи в разработанной системе реализован в соответствии со стандартом PKCS\#7~\cite{pkcs7}, который определяет синтаксис криптографических сообщений. Современная версия этого стандарта известна как Cryptographic Message Syntax (CMS)~\cite{rfc5652} и описана в RFC 5652. PKCS\#7/CMS определяет структуру данных для цифровых подписей, включая информацию о подписанте, используемых алгоритмах, сертификатах и самой подписи.


Стандарт PKCS\#1~\cite{rfc3447}, описанный в RFC 3447, определяет спецификации криптографии с открытым ключом на основе RSA, включая форматы ключей и схемы подписи. В разрабатываемой системе используется схема подписи RSASSA-PKCS1-v1\_5, широко поддерживаемая криптографическими библиотеками.

\subsection{Нормативно-правовая база}

Федеральный закон № 63-ФЗ «Об электронной подписи»~\cite{fz63} устанавливает правовые условия использования электронных подписей в Российской Федерации. Закон определяет три вида электронных подписей и условия их применения. Для внутрикорпоративного документооборота, который является целью данной работы, подходит усиленная неквалифицированная электронная подпись (НЭП), которая создается с использованием криптографических средств и позволяет идентифицировать подписанта и обнаружить факт изменения документа после подписания. При этом разрабатываемая система обеспечивает поддержку отсоединенной КЭП, что расширяет её применимость для различных сценариев документооборота.

Статья 22.3 Трудового кодекса РФ~\cite{tk_rf} регламентирует использование электронного документооборота в трудовых отношениях. Согласно этой статье, работодатель и работник могут использовать электронные документы при условии соблюдения требований к их созданию, подписанию и хранению. Система обеспечивает техническую реализацию этих требований.

\subsection{Существующие программные решения}

Существует несколько категорий решений для работы с электронной подписью:

\textbf{Коммерческие СКЗИ.} Сертифицированные средства криптографической защиты информации (например, КриптоПро CSP, ViPNet CSP) предоставляют полный набор функций для работы с квалифицированной электронной подписью. Однако их использование связано со значительными затратами на лицензирование, требует интеграции со специфичными API и ориентировано на работу с аккредитованными удостоверяющими центрами. Для внутрикорпоративного документооборота, не требующего квалифицированной подписи, применение таких решений экономически нецелесообразно.

\textbf{Облачные сервисы ЭЦП.} Сервисы типа DocuSign, Adobe Sign предоставляют возможности электронного подписания через облачную инфраструктуру. Недостатками таких решений являются зависимость от внешнего провайдера, ограничения в кастомизации процессов, необходимость передачи конфиденциальных документов на сторонние серверы и регулярные платежи за использование сервиса.

\textbf{Библиотеки с открытым исходным кодом.} OpenSSL~\cite{openssl} --- широко используемая криптографическая библиотека, предоставляющая низкоуровневые примитивы для работы с RSA, хешированием и другими криптографическими операциями. Python-библиотека cryptography~\cite{cryptography_python} предоставляет высокоуровневый интерфейс к криптографическим операциям. Библиотека pyHanko~\cite{pyhanko} специализируется на подписании PDF-документов. Эти инструменты предоставляют строительные блоки для создания системы электронной подписи, но требуют разработки полноценного решения на их основе.

\subsection{Обоснование выбранного подхода}

Анализ существующих решений показывает, что для задачи обеспечения внутрикорпоративного документооборота наиболее целесообразным являлась разработка собственного решения на основе открытых криптографических библиотек и стандартов. Такой подход позволяет:

\begin{enumerate}
    \item Избежать зависимости от коммерческих СКЗИ и связанных с ними затрат, сохраняя при этом криптографическую стойкость за счет использования проверенных алгоритмов (RSA-2048, SHA-256) и стандартов (PKCS\#7).

    \item Обеспечить полный контроль над процессами генерации ключей, подписания и хранения данных, что важно для безопасности и соответствия внутренним политикам организации.

    \item Создать API, максимально интегрированный с существующей корпоративной платформой, в отличие от универсальных облачных сервисов.

    \item Гарантировать соответствие требованиям ФЗ-63~\cite{fz63} для усиленной неквалифицированной электронной подписи.
\end{enumerate}

Выбор микросервисной архитектуры с разделением C++ бэкенда (userver~\cite{userver}) и Python криптографического сервиса обусловлен следующими соображениями:

\begin{itemize}
    \item C++ обеспечивает высокую производительность для обработки большого количества HTTP-запросов и работы с базой данных~\cite{postgresql}.

    \item Python предоставляет богатую экосистему криптографических библиотек (cryptography~\cite{cryptography_python}, pyHanko~\cite{pyhanko}), упрощающих реализацию сложных криптографических операций.

    \item Разделение на микросервисы позволяет независимо масштабировать компоненты и упрощает обновление криптографической логики без затрагивания основного бэкенда.
\end{itemize}

\subsection{Позиционирование работы}

Данная работа создает оптимизированное решение для внутрикорпоративного документооборота, сочетающее криптографическую стойкость стандартных алгоритмов, экономическую эффективность открытых технологий и гибкость интеграции с существующей платформой. В отличие от универсальных коммерческих решений, разрабатываемая система не перегружена избыточной функциональностью и оптимизирована для конкретных требований организации.

\section{Реализация: архитектура, данные и публичный API}

\subsection{Назначение и границы подсистемы}

Подсистема НЭП в составе платформы обеспечивает: (1)~однократную генерацию пары RSA-2048 для сотрудника, (2)~подписание PDF, уже размещённого в объектном хранилище, с формированием отсоединённой подписи PKCS\#7 (файл \ppath{.p7s}), (3)~проверку сохранённой подписи с раздельными признаками целостности файла и криптографической корректности. Публичный контур НЭП не обращается к удостоверяющему центру; для НЭП в CMS используется самоподписанный сертификат.

Один и тот же \ppath{document_id} может иметь \textbf{несколько} записей в \ppath{document_signatures} (разные сотрудники, разные моменты подписания): каждая подпись хранится отдельным объектом в S3 и проверяется по своему \ppath{signature_id}. Параллельно платформа поддерживает \textbf{КЭП}, сформированную на рабочем месте пользователя: во фронтенд встроен плагин для работы с квалифицированной подписью, после чего отсоединённый файл подписи загружается на бэкенд (маршрут \httpapi{POST /v1/documents/upload-signature}, тип \ppath{kep} в метаданных). Сочетание нескольких НЭП и одной или нескольких КЭП для одного документа не отменяет криптографическую валидность: каждая подпись независимо проверяется относительно того состояния файла, которое было захешировано при её создании (при неизменном PDF все подписи остаются корректными).

\subsection{Логическая архитектура}

\textbf{Компоненты}

Userver (C++) - единая точка входа для внешних клиентов: аутентификация, авторизация, работа с БД компании, оркестрация вызова Python-сервиса.

Python\_service (aiohttp) - построение CMS для PDF, чтение и запись байтов документа и подписи через общий S3-клиент.

PostgreSQL - ключи сотрудника, метаданные подписей, пароль подписи (отдельная сущность), журналы аудита (см. миграции).

S3-совместимое хранилище - тело PDF по ключу, согласованному с \ppath{document_id}, файл \ppath{.p7s} по пути из метаданных.

Регистрация маршрутов и URL вызовов Python-сервиса задаются в \ppath{configs/static_config.yaml}. Имена компонентов: \ppath{handler-v1-employee-keys-generate}, \ppath{handler-v1-documents-nep-sign}, \ppath{handler-v1-documents-nep-verify}.
У генерации ключей нет \ppath{pyservice-url} (только C++/OpenSSL). У \ppath{nep-sign} и \ppath{nep-verify} полный путь совпадает с маршрутами aiohttp в \ppath{python_service/main.py}.

Реализации на C++: \ppath{src/views/v1/employee/keys/generate/view.cpp}, \ppath{src/views/v1/documents/nep_sign/view.cpp}, \ppath{src/views/v1/documents/nep_verify/view.cpp}.

Внутренние маршруты Python-приложения: \httpapi{POST /document/nep-sign}, \httpapi{POST /document/nep-verify} (\ppath{python_service/main.py}).

Байты PDF и \ppath{.p7s} не проксируются через C++: загрузка и выгрузка выполняется в Python, что снижает нагрузку на сеть. При подписании приватный PEM передаётся во внутренний вызов уже после расшифровки на стороне C++.

\begin{table}[htbp]
\centering
\caption{Участие компонентов в операциях НЭП}
\small
\begin{tabular}{p{2.4cm}p{3.1cm}p{2.8cm}p{2.2cm}p{2.4cm}}
\toprule
Операция & C++ & Python & S3 & PostgreSQL \\
\midrule
Генерация ключей & RSA, PEM, хеш & --- & --- & \ppath{employee_keys}, пароль \\
\midrule
Подписание & проверки, расшифровка PEM, HTTP в Py & CMS, \ppath{.p7s} & чтение PDF, запись подписи & \ppath{document_signatures} \\
\midrule
Проверка & метаданные, HTTP в Py & проверка CMS & чтение PDF и \ppath{.p7s} & чтение по \ppath{signature_id} \\
\bottomrule
\end{tabular}
\end{table}



Клиент обращается к единому HTTP API на C++ (userver~\cite{userver}); идентификаторы пользователя и компании извлекаются из контекста Bearer-сессии. Операции подписания и проверки инициируют HTTP-запросы к внутреннему \ppath{python_service} (aiohttp, порты и URL задаются в \ppath{configs/static_config.yaml}). Тот читает и пишет байты PDF и файлов \ppath{.p7s} в S3-совместимое хранилище; метаданные и ключи фиксируются в PostgreSQL~\cite{postgresql} в схеме компании.

Схема БД вводится миграцией \ppath{26__nep_signature.sql}.
\begin{verbatim}
CREATE TABLE IF NOT EXISTS ${SCHEMA}.employee_keys (
    employee_id TEXT PRIMARY KEY NOT NULL,
    private_key TEXT NOT NULL,
    public_key TEXT NOT NULL,
    ...
);
CREATE TABLE IF NOT EXISTS ${SCHEMA}.document_signatures (
    id TEXT PRIMARY KEY NOT NULL,
    document_id TEXT NOT NULL,
    employee_id TEXT NOT NULL,
    signature_path TEXT NOT NULL,
    signature_metadata JSONB NOT NULL,
    ...
);
\end{verbatim}
{\small В файле миграции имя схемы подставляется параметром для каждой компании.}

Дополнительно вводятся связанные объекты: хранение пароля подписи - \ppath{27__employee_signature_password.sql}; поле типа подписи (nep/kep) - \ppath{28__signature_type.sql}; аудит изменений пароля подписи - \ppath{34__employee_signature_passwords_audit.sql}. Актуальное состояние схемы компании следует сверять с каталогом \ppath{postgresql/migrations/}.

Публичный контракт трёх маршрутов НЭП зафиксирован в машиночитаемой спецификации \ppath{docs/yaml/openapi.yaml}. Краткое содержание контракта:

\httpapi{POST /v1/employee/keys/generate} --- query \ppath{signature_password} (ровно 6 цифр). Ответ 200: \ppath{public_key} (PEM), \ppath{public_key_hash}. Типовая ошибка 400: ключи для сотрудника уже существуют.

\httpapi{POST /v1/documents/nep-sign} --- query \ppath{document_id}, \ppath{signature_password}; тело JSON опционально: \ppath{reason}, \ppath{location}. Ответ 200: \ppath{signature_id}, \ppath{signature_path}, \ppath{timestamp}, \ppath{public_key_hash}. Коды 400 (нет ключей), 401 (неверный пароль), 404 (документ), 502 (S3 или внутренний сервис).

\httpapi{GET /v1/documents/nep-verify} --- query \ppath{signature_id}. Ответ 200: \ppath{valid}, \ppath{integrity_ok}, \ppath{signature_ok}, \ppath{message}, \ppath{verified_at}; опционально \ppath{signer_name}, \ppath{signature_timestamp}. Коды 404, 502 --- по смыслу отсутствия записи или сбоя чтения.

Внутренний HTTP API Python-сервиса вызывается только из доверенного контура (C++ $\to$ URL из \ppath{configs/static_config.yaml}) и не предназначен для прямого вызова внешними клиентами.

\httpapi{POST /document/nep-sign}: тело JSON с обязательными полями \ppath{document_id}, \ppath{employee_id}, \ppath{employee_name}, \ppath{private_key}, \ppath{public_key} (строки PEM) и опционально \ppath{reason}, \ppath{location}. Поведение: загрузка PDF из S3 по ключу, совпадающему с \ppath{document_id} (как ожидает \ppath{python_service/s3_client}), формирование отсоединённой подписи, запись файла подписи в bucket под именем вида
\ppath|{document_id}_{employee_id}_nep.p7s|. Путь к объекту подписи в S3 сохраняется в \ppath{document_signatures.signature_path} и используется при последующей проверке.
Успешный ответ JSON содержит \ppath{signature_path}, \ppath{timestamp}, \ppath{public_key_hash}, \ppath{user_name}, \ppath{user_id} и при наличии --- \ppath{reason}, \ppath{location}. Идентификатор \ppath{signature_id} назначается на стороне C++ при записи строки в \ppath{document_signatures}.

\httpapi{POST /document/nep-verify}: обязательные \ppath{document_id}, \ppath{signature_path}, \ppath{public_key}; загрузка PDF и \ppath{.p7s}, вызов \ppath{NEPSigner.verify_pdf}. В ответ передаётся набор полей, который после обработки на стороне C++ формирует тело ответа \httpapi{GET /v1/documents/nep-verify} (\ppath{valid}, \ppath{integrity_ok}, \ppath{signature_ok}, \ppath{message}, \ppath{verified_at}; при наличии --- \ppath{signer_name}, \ppath{signature_timestamp}).

\httpapi{POST /document/build-archive}: внутренний маршрут для сборки ZIP. Тело JSON: \ppath{document_id}, массив \ppath{signature_paths}. Реализация --- \ppath{python_service/build_archive/handler.py}: скачивание тела документа и каждой подписи из S3, упаковка в \ppath{.zip}, загрузка архива в S3, возврат presigned URL.

\subsection{Цепочка согласования, КЭП и выгрузка архива}

\textbf{Цепочка подписания (approval chain).} Для поэтапного согласования документа используются маршруты \httpapi{POST /v1/documents/chain/add} (привязка к документу описания цепочки: кто и в каком порядке действует) и \httpapi{POST /v1/documents/chain/update?document_id=...} (переход текущего шага). В метаданных шага задаётся, требуется ли подпись и какого типа. Для шага с НЭП (\ppath{requires_signature = 2}) пользователь передаёт \ppath{signature_password}; бэкенд расшифровывает PEM и вызывает тот же внутренний \ppath{/document/nep-sign}, что и прямой маршрут \ppath{/v1/documents/nep-sign}, после чего в \ppath{document_signatures} добавляется запись с \ppath{signature_type = nep}. Для шага с КЭП (\ppath{requires_signature = 1}) пользователь сначала получает подпись на клиенте (плагин КЭП), загружает файл через \ppath{/v1/documents/upload-signature}, затем на шаге chain передаёт \ppath{signature_id} созданной записи --- бэкенд проверяет, что подпись принадлежит текущему пользователю, документу и типу \ppath{kep}, и переводит шаг цепочки в состояние «подписано». Таким образом, несколько человек последовательно фиксируют согласие, не перезаписывая чужие подписи.

\textbf{Скачивание документа со всеми подписями.} Публичный маршрут \httpapi{GET /v1/documents/download-with-signatures?id=<document_id>} (обработчик \ppath{src/views/v1/documents/download_with_signatures/view.cpp}) после проверки доступа читает из БД все \ppath{signature_path} для данного \ppath{document_id} и вызывает Python \httpapi{POST /document/build-archive} (в \ppath{static_config.yaml} задаётся как \ppath{pyservice-url} для этого обработчика). Клиент получает JSON с presigned URL; по ссылке скачивается \textbf{ZIP-архив}, содержащий файл документа и все отсоединённые подписи (\ppath{.p7s} и т.\,п.) --- удобно для архива, передачи контрагенту и внешней проверки каждой подписи средствами, понимающими PKCS\#7.

\subsection{Система с точки зрения пользователя}

Пользователь платформы (после обычного входа и получения Bearer-токена) выполняет типовой цикл НЭП: (1)~однократно задаёт шестизначный пароль подписи и вызывает генерацию ключей; (2)~для подписания PDF передаёт \ppath{document_id} и пароль, при необходимости --- причину и место подписания в JSON; (3)~для проверки передаёт \ppath{signature_id} и получает признаки \ppath{valid}, \ppath{integrity_ok}, \ppath{signature_ok} и текст \ppath{message}. Тот же пользователь (или другие участники цепочки) могут добавить ещё подписи к тому же документу; при согласовании по цепочке шаг с НЭП или КЭП выполняется в интерфейсе согласования (см. выше). Для КЭП без серверных ключей пользователь подписывает документ плагином в браузере и загружает файл подписи. Чтобы получить один файл для хранения или передачи, используют \httpapi{GET /v1/documents/download-with-signatures} --- в ответ приходит ссылка на ZIP с документом и всеми подписями. Поля ответов и коды ошибок согласованы с OpenAPI.

Проверка целостности (\ppath{integrity_ok}) отражает совпадение хеша документа с данными подписи; \ppath{signature_ok} --- криптографическую корректность подписи относительно переданного публичного ключа; \ppath{valid} логически соответствует конъюнкции этих условий.

Примеры вызовов без отдельного клиентского приложения:

\begin{verbatim}
curl -X POST "https://<host>/v1/employee/keys/generate?signature_password=123456" \
  -H "Authorization: Bearer <token>"

curl -X POST "https://<host>/v1/documents/nep-sign?document_id=<uuid>&signature_password=123456" \
  -H "Authorization: Bearer <token>" -H "Content-Type: application/json" \
  -d '{"reason":"Согласование","location":"Москва"}'

curl -X GET "https://<host>/v1/documents/nep-verify?signature_id=<uuid>" \
  -H "Authorization: Bearer <token>"
\end{verbatim}

\subsection{Типовые отказы и действия}

\textbf{Сообщение о ненайденных ключах (код 400).} Ключевая пара для сотрудника не создавалась: выполнить \httpapi{POST /v1/employee/keys/generate} с корректным \ppath{signature_password}.

\textbf{Недействительная подпись или \ppath{integrity_ok=false}.} Возможные причины: изменение PDF после подписания, повреждение файла \ppath{.p7s}, несоответствие публичного ключа. Действия: восстановить исходный документ при необходимости, повторить подписание.

\textbf{Сообщения об ошибке целостности PDF (в т.\,ч. повреждённый файл).} Восстановить корректный PDF, затем подписать заново.

\subsection{Объёмные характеристики реализации}

Оценка относится только к подсистеме НЭП, без объёма всей платформы. Числа получены подсчётом строк в репозитории на момент подготовки отчёта.

\begin{table}[htbp]
\centering
\caption{Объём исходного кода и тестов подсистемы НЭП}
\begin{tabular}{lrr}
\toprule
Компонент & Строк (прибл.) & Объём, Кбайт \\
\midrule
C++: три обработчика (\ppath{keys/generate}, \ppath{nep_sign}, \ppath{nep_verify}) & 732 & $\approx$26 \\
Python: пакет \ppath{nep_signer} & 706 & $\approx$31 \\
Автотесты: \ppath{tests/test_nep_basic.py}, \ppath{python_service/tests/test_nep_signer.py}, \ppath{python_service/tests/test_nep_handlers.py} & 1325 & $\approx$43 \\
SQL: миграция \ppath{26__nep_signature.sql} & $\approx$37 & $\approx$2{,}4 \\
\bottomrule
\end{tabular}
\end{table}

В миграции задаются две таблицы и три индекса на схему компании; физический объём БД определяется числом сотрудников и подписей и в отчёте не нормируется.

\subsection{Сборка, зависимости и применение миграций}

\textbf{Python.} Зависимости перечислены в \ppath{python_service/requirements.txt}: в первую очередь \texttt{cryptography}, \texttt{pyhanko}, \texttt{asn1crypto} (разбор структур PKCS\#7), а также пакеты для aiohttp и клиента S3; для разработки --- \ppath{requirements-dev.txt}. Установка в виртуальное окружение и прогон юнит-тестов из каталога \ppath{python_service}: \ppath|pip install -r requirements-dev.txt|, \ppath|python3 -m pytest python_service/tests/ -v|; из корня репозитория допускается цель \ppath{make test-nep-unit} (при необходимости создаётся \ppath{python_service/.venv}).

\textbf{C++.} Обработчики подключаются в сборку CMake основного проекта платформы стандартным для репозитория способом.

\textbf{Миграции PostgreSQL.} Файлы в \ppath{postgresql/migrations/} применяют к схеме каждой компании согласно принятой в проекте процедуре; пример одной миграции НЭП приведён в приложении~Б.

\textbf{Выводы раздела.} Выбрана микросервисная декомпозиция «оркестратор + криптосервис + объектное хранилище»; границы ответственности и форматы полей согласованы между OpenAPI, SQL и внутренним JSON Python-обработчиков. Порядок величин объёма кода и тестов фиксирует масштаб доработки при интеграции в существующий продукт.

\section{Реализация: криптографические модули}

\textbf{C++, OpenSSL.} Генерация RSA-2048 и запись приватного ключа в PEM с шифрованием AES-256-CBC от пароля подписи реализованы в \ppath{src/views/v1/employee/keys/generate/view.cpp}. Существенный фрагмент:

\begin{lstlisting}[language=C++, basicstyle=\small\ttfamily, columns=fullflexible]
if (EVP_PKEY_CTX_set_rsa_keygen_bits(ctx, 2048) <= 0) { ... }
if (PEM_write_bio_PrivateKey(private_bio, pkey,
        EVP_aes_256_cbc(), kstr, klen, nullptr, nullptr) != 1) { ... }
SHA256(reinterpret_cast<const unsigned char*>(result.public_key.c_str()),
       result.public_key.length(), hash);
\end{lstlisting}

Расшифровка PEM при подписании выполняется в обработчике \ppath{nep_sign} до вызова Python; при неверном пароле цепочка обрывается с кодом 401 (см. OpenAPI).

\textbf{Python, класс \texttt{NEPSigner}.} Файл \ppath{python_service/nep_signer/signer.py}: загрузка пары ключей из PEM, метод \ppath{sign_pdf} (отсоединённая CMS/PKCS\#7 для PDF, самоподписанный сертификат, SHA-256), метод \ppath{verify_pdf} с полями \ppath{integrity_ok} и \ppath{signature_ok}. HTTP-обёртки \ppath{sign_handler.py} и \ppath{verify_handler.py} скачивают объекты из S3 по \ppath{document_id} и \ppath{signature_path}.

\textbf{Выводы раздела.} Криптографическая цепочка соответствует требованиям введения: RSA-2048, SHA-256, PKCS\#7/CMS~\cite{pkcs7,rfc5652,rfc3447}; разделение «генерация/хранение ключа на C++ --- CMS на Python» снижает объём криптокода в основном бэкенде.

\section{Тестирование и сопоставление с аналогами}

\subsection{Метрики качества и обоснование достоверности}

В смысле п.~2.5 методических рекомендаций заявления о корректности подкрепляются не рассуждением, а \emph{воспроизводимыми} проверками. Выбранные метрики: (1)~покрытие основных и альтернативных веток API (коды 400/401/404/502); (2)~изоляция криптографии в юнит-тестах без сети; (3)~наличие негативных сценариев (искажённый PDF, неверный пароль, отказ pyservice/S3).

Суммарно в трёх файлах автотестов реализовано \textbf{57} именованных тестовых сценариев (19 интеграционных через userver с моками Python-сервиса, 24 на \texttt{NEPSigner}, 14 на HTTP-обработчики Python). Время полного прогона зависит от машины и CI и отдельно не протоколировалось; порядок --- единицы минут при локальном запуске подмножества. Нагрузочный эксперимент (устойчивый RPS) \textbf{не проводился} --- это ограничение полноты экспериментальной части и явно учтено в заключении.

\subsection{Состав автотестов}

\ppath{tests/test_nep_basic.py} --- генерация ключей, \ppath{nep-sign}, \ppath{nep-verify}, ошибки пароля, отсутствие документа/ключей, ответы pyservice 500/502, отклонение не-НЭП подписи. \ppath{python_service/tests/test_nep_signer.py} --- цикл подпись--проверка, порча PDF и \ppath{.p7s}, ветки \ppath{integrity_ok}/\ppath{signature_ok}, размеры PDF (параметризация). \ppath{python_service/tests/test_nep_handlers.py} --- обязательные поля JSON, ошибки S3, успех и ошибки верификации.

\subsection{Производительность и сопоставление с аналогами}

Оценка вычислительной сложности на уровне продукта ограничена: доминируют I/O (S3, HTTP между сервисами) и криптопримитивы библиотек; асимптотический анализ не выносился, так как он не добавляет ценности поверх стандартных RSA/SHA-256.

\subsubsection{Скрипт замера задержек публичных обработчиков}

Замеряется именно время ответа HTTP-ручек userver (то, что видит клиент); внутри \ppath{nep-sign} и \ppath{nep-verify} в него входят обращения к Python-сервису и S3. Скрипт \ppath{scripts/bench_nep_handlers.sh} выполняет последовательно: (1)~вставку в PostgreSQL временного сотрудника, строки в \ppath{wd_general.auth_tokens} и документа в схеме компании (в тестовой конфигурации используется компания с идентификатором \ppath{first}); (2)~при заданном bucket --- генерацию минимального PDF и загрузку в S3 с ключом, равным \ppath{document_id}; (3)~серии запросов \ppath{curl} и вывод статистики min/max/mean для \ppath{time_total} (секунды или миллисекунды в зависимости от настроек \ppath{curl}); для \httpapi{POST /v1/employee/keys/generate} число повторов задаётся \ppath{NEP_BENCH_KEYS_ITERATIONS}, между итерациями скрипт удаляет строки \ppath{employee_keys} и \ppath{employee_signature_passwords} у бенч-сотрудника (без этого повторный запрос вернул бы 400) --- \emph{время этих DELETE в статистику не входит}; при наличии S3 аналогично замеряются \ppath{nep-sign} и \ppath{nep-verify} с числом повторов \ppath{NEP_BENCH_ITERATIONS}; (4)~по завершении и при ошибке (\ppath{trap}) удаление вставленных данных для чистоты повторного запуска.

Требования к окружению: установленные \ppath{psql}, \ppath{curl}, \ppath{jq}; развёрнутое виртуальное окружение \ppath{python_service/.venv} с зависимостями; запущенный API; доступ к PostgreSQL. Токен Bearer скрипт не запрашивает у оператора --- используется запись из \ppath{auth_tokens}, вставленная скриптом.

\begin{table}[htbp]
\centering
\caption{Переменные окружения сценария \ppath{scripts/bench_nep_handlers.sh}}
\small
\begin{tabular}{p{4.5cm}p{8.5cm}}
\toprule
Переменная & Назначение \\
\midrule
\ppath{NEP_BENCH_BASE_URL} & Базовый URL API (по умолчанию \ppath{http://127.0.0.1:8080}) \\
\ppath{NEP_BENCH_ITERATIONS} & Число повторов для \ppath{nep-sign} и \ppath{nep-verify} (по умолчанию 7) \\
\ppath{NEP_BENCH_KEYS_ITERATIONS} & Повторов для \ppath{keys/generate} (по умолчанию то же; между вызовами --- очистка ключей в БД) \\
\ppath{NEP_BENCH_SIGNATURE_PASSWORD} & Пароль подписи для generate/\ppath{nep-sign} (по умолчанию \texttt{123456}) \\
\ppath{NEP_BENCH_PGURI} & URI PostgreSQL \emph{или} стандартные \ppath{PGHOST}, \ppath{PGPORT}, \ppath{PGUSER}, \ppath{PGPASSWORD}, \ppath{PGDATABASE} \\
\ppath{NEP_BENCH_S3_BUCKET} & Если не задан --- выполняется только замер \ppath{keys/generate} \\
\ppath{AWS_ACCESS_KEY_ID}, \ppath{AWS_SECRET_ACCESS_KEY} & Учётные данные для загрузки в S3 \\
\ppath{AWS_ENDPOINT_URL} или \ppath{NEP_BENCH_S3_ENDPOINT} & Endpoint S3-совместимого хранилища при необходимости \\
\ppath{AWS_DEFAULT_REGION} & Регион (в вспомогательных сценариях загрузки по умолчанию может быть \ppath{ru-central1}) \\
\bottomrule
\end{tabular}
\end{table}

\textbf{Пример результатов одного прогона} (\(n=7\) повторов на каждую ручку при значениях по умолчанию скрипта; одна конфигурация машины и окружения разработки, не нагрузочный тест и не статистически репрезентативная выборка):

\begin{center}
\small
\begin{tabular}{lrrr}
\toprule
Эндпоинт & min, мс & max, мс & mean, мс \\
\midrule
\httpapi{POST /v1/employee/keys/generate} & 124{,}4 & 344{,}5 & 213{,}5 \\
\httpapi{POST /v1/documents/nep-sign} & 112{,}5 & 327{,}9 & 172{,}6 \\
\httpapi{GET /v1/documents/nep-verify} & 45{,}4 & 123{,}0 & 84{,}1 \\
\bottomrule
\end{tabular}
\end{center}

\textbf{Вывод по этим цифрам} (качественный, без претензии на норму для продакшена): средняя задержка \ppath{nep-verify} примерно вдвое ниже, чем у \ppath{nep-sign} и \ppath{keys/generate}, что согласуется с тем, что верификация не порождает новую подпись и в типичном сценарии меньше по объёму работы в Python/S3, тогда как подпись и генерация пары включают тяжёлые криптооперации и обмен с вспомогательным сервисом. Широкий интервал min--max при \(n=7\) ожидаем: в него попадают прогрев, конкуренция за CPU/диск, сеть до S3 и БД; для сравнения с внешними продуктами или для SLA нужны отдельные серии замеров на контролируемом стенде.

Между итерациями \ppath{keys/generate} выполнялись DELETE в БД (в замер не входят), как описано выше.

Запуск из корня репозитория: \ppath{make bench-nep} или \ppath|bash scripts/bench_nep_handlers.sh| после экспорта переменных (URI БД, при необходимости --- bucket и ключи S3).

\subsubsection{Ручной разовый замер и границы измерения}

Для одного запроса к уже существующему пользователю можно использовать:

\begin{verbatim}
curl -s -o /dev/null -w "time_total=%{time_total}s\n" \
  -X POST "${NEP_BENCH_BASE_URL}/v1/employee/keys/generate?signature_password=123456" \
  -H "Authorization: Bearer YOUR_TOKEN"
\end{verbatim}

Время выполнения юнит-тестов \ppath{python_service/tests/test_nep_signer.py} отражает работу библиотеки в процессе Python и \textbf{не} эквивалентно задержке публичных ручек userver.

\subsubsection{Ручная приёмка и мониторинг}

Рекомендуемый ручной контроль перед сдачей: сгенерировать ключи тестовому пользователю, подписать PDF, выполнить проверку, изменить PDF и убедиться, что проверка фиксирует нарушение целостности. В эксплуатации целесообразно по журналам отслеживать объёмы операций генерации ключей и подписания, а также ошибки подписи и верификации.

Сопоставление с аналогами (функциональность / эффективность / удобство): коммерческие СКЗИ дают сертифицированный контур КЭП и УЦ, но выше TCO и жёстче интеграция; облачные ЭЦП снижают нагрузку на ИТ, но выносят документы к провайдеру. Предложенная система выигрывает в \textbf{контроле данных} и \textbf{стоимости владения} для НЭП-внутряка, уступая в сертификации СКЗИ и в «готовности из коробки» по сравнению с вендором.

\textbf{Выводы раздела.} Достоверность и полнота в рамках выбранных метрик обеспечиваются 57 автотестами и описанием метрик в настоящем разделе; отрицательный результат --- отсутствие нагрузочного отчёта. Сравнение с рынком аргументирует практическую нишу решения.

\section{Заключение}

Сопоставление с критериями оценки результатов (актуальность, новизна, теоретическая и практическая значимость, достоверность, полнота):

\begin{itemize}
    \item \textbf{Актуальность и практическая полезность.} Внутренний ЭДО и ст.~22.3 ТК~РФ~\cite{tk_rf}, ФЗ~63~\cite{fz63} задают спрос на НЭП без УЦ; реализация встраивается в действующую платформу (разд.~3--4).
    \item \textbf{Новизна.} На уровне алгоритмов новизны нет; на уровне системной интеграции --- связка userver + выделенный криптосервис + PKCS\#7 + раздельная верификация (разд.~3--5).
    \item \textbf{Теоретическая опора.} Обзор литературы и стандартов RSA, SHA-256, CMS (разд.~2).
    \item \textbf{Достоверность и корректность.} Подтверждаются 57 автотестами и ссылками на файлы кода; криптографические параметры и потоки данных согласованы с разд.~3--4 и постановкой задачи (разд.~1).
    \item \textbf{Полнота.} Ограничена отсутствием нагрузочного эксперимента и формального отчёта по модели угроз для секретов в БД (отрицательные результаты/пробелы зафиксированы явно).
\end{itemize}

\textbf{Положительные результаты:} работающий контур НЭП (три основных эндпоинта) и интеграция с платформенными сценариями: несколько подписей на один документ, КЭП с клиента (плагин + \ppath{upload-signature}), цепочка согласования (\ppath{chain/add}, \ppath{chain/update}), выгрузка ZIP через \ppath{download-with-signatures}; миграция БД и сопутствующие объекты (пароль подписи, тип подписи, аудит); $\approx$730 строк C++ и $\approx$710 строк Python только по НЭП; машиночитаемый контракт API (\ppath{docs/yaml/openapi.yaml}); в отчёте изложены назначение и архитектура НЭП, схема данных, публичный и внутренний API (в т.\,ч. фрагмент \ppath{static_config.yaml}), соглашения S3, нормативная база и обоснование подхода в разд.~2, пользовательский сценарий, типовые отказы, сборка и миграции, объёмы кода, тестирование (57 сценариев), протокол и переменные скрипта замера задержек, ручная приёмка и мониторинг; в приложении~Б --- пример команды применения миграции; возможность демонстрации через развёрнутую платформу и автотесты.

\textbf{Отрицательные результаты / пробелы:} нет измеренного RPS и профиля задержек под нагрузкой; нет сертификации СКЗИ; анализ хранения ключей и пароля в БД --- вынесен в перспективу.

\textbf{Перспективы:} нагрузочные испытания; развитие административных инструкций; при необходимости --- отдельный документ по угрозам и политике секретов.

\newpage
\printbibliography[heading=bibintoc]

\appendix
\section{Сокращения и обозначения}

\begin{description}
    \item[НЭП] усиленная неквалифицированная электронная подпись.
    \item[КЭП] квалифицированная электронная подпись.
    \item[CMS] Cryptographic Message Syntax (RFC~5652), развитие PKCS\#7.
    \item[PEM] текстовое кодирование криптоключей (Privacy-Enhanced Mail).
    \item[REST] архитектурный стиль HTTP API.
    \item[S3] API объектного хранилища (в проекте --- совместимое с S3).
    \item[userver] фреймворк асинхронного сервера на C++ для сервисов платформы.
\end{description}

\section{Пример применения миграции схемы НЭП}

Миграции применяют к схеме каждой компании в порядке номеров файлов; ниже приведён пример для одной схемы (подставить имя базы, пользователя СУБД и идентификатор компании в имя схемы):

\begin{verbatim}
psql -U postgres -d working_day -v SCHEMA=working_day_<company_id> \
  -f postgresql/migrations/26__nep_signature.sql
\end{verbatim}

Далее для той же схемы последовательно применяют \ppath{27__employee_signature_password.sql}, \ppath{28__signature_type.sql} и остальные файлы из \ppath{postgresql/migrations/}, относящиеся к подписи и документам, согласно регламенту развёртывания проекта.

\end{document}
